\begin{lemma}{Elementare Eigenschaften der Gamma-- und Betafunktion}
             {ElementareEigenschaftenDerGammafunktion}
Seien $\Gamma:\R_{>0}\rightarrow\R,\,
t\mapsto\int_0^\infty x^{t-1}\exp(-x)dx$ und 
\[B: (\R_{>0})^2 \rightarrow \R,
(x,y)\mapsto\int_0^1 (1-s)^{x-1}s^{y-1}ds\text.\]
Dann gilt:
\begin{enumerate}[(1)]
	\item \label{ElementareEigenschaftenDerGammafunktion1}
	$\forall\, t \in \R_{>0}: \Gamma(t) > 0$.
	\item \label{ElementareEigenschaftenDerGammafunktion2}
	$\forall\, t \in \R_{>0}: \Gamma(t+1) = t \Gamma(t)$.
	\item\label{ElementareEigenschaftenDerGammafunktion3}
	$\forall\, n \in \N: \Gamma(n) = (n-1)!$
	\item\label{ElementareEigenschaftenDerGammafunktion4}
	$\forall\,x,y\in\R_{>0}: B(x,y) = \frac{\Gamma(x)\Gamma(y)}{\Gamma(x+y)}$.
\end{enumerate}
\end{lemma}
\begin{proof}
(1) Sei $t \in \R_{>0}$.\\
Sei $f_t: \R_{>0} \rightarrow \R, x \mapsto x^{t-1}\exp(-x)$.\\
Dann gilt f�r alle
$x \in \R_{>0}: f_t(x) = \exp((t-1)\log(x))\exp(-x) > 0$, also $f_t > 0$.\\
Da $f_t$ offensichtlich stetig ist, ist auch
$\Gamma(t)=\int_0^\infty f_t(y)dy > 0$.\\
(2)/(3) siehe \cite{K1} dort 18.1.(5) und 18.1.(4).\\
(4) siehe \cite{K2}, dort 8.4.(8).
\end{proof}